\section{Actual orbits and sharp seed counts}\label{supp3:section:001}
Let $C\in M_{d\times m}(\C)$ be the coefficient matrix of one seed. Define its orbit by actual matrix multiplication:
\[
\mathcal O(C)=\{AC:A\in M_d(\C)\}.
\]
This set is already a linear subspace. The source proves that it equals the linear span of the orbit before computing its dimension.

A matrix $X$ belongs to $\mathcal O(C)$ exactly when each of its rows belongs to $\operatorname{Row}(C)$. One direction follows by expanding $AC$. For the converse, choose the coefficients producing each row separately and place them in the corresponding row of $A$. An explicit linear equivalence therefore gives
\[
\mathcal O(C)\simeq \operatorname{Row}(C)^d,
\qquad
\boxed{\dim_\C\mathcal O(C)=d\,\rank C.}
\]
If $d>0$, equality with the full $d\times m$ coefficient space is equivalent to $\rank C=m$. Consequently a cyclic seed requires $m\le d$. The assumption that there is at least one row is essential for canceling the factor $d$.

\subsection{Several seeds: both necessity and construction}
Continue under the assumption above that the number of rows is positive. For seeds $C_1,\ldots,C_k$, generation means the concrete statement
\[
\forall X\in M_{d\times m}(\C),\quad
\exists A_1,\ldots,A_k\in M_d(\C):\quad
X=\sum_{a=1}^k A_aC_a.
\]
Stack the $C_a$ vertically into a $kd\times m$ matrix $D$. The source identifies the action on the stacked matrix and proves the equivalence
\[
\text{the seeds generate}\quad\Longleftrightarrow\quad\rank D=m.
\]
Thus $m\le kd$ is necessary. It is sufficient as well: embed the $m$ coordinate columns into $kd$ coordinate positions to obtain a matrix $D$ of full column rank, then split its rows back into $k$ seeds. This construction proves that generating matrices with the required dimensions exist.

\subsection{Simultaneous generation in a direct sum}
For the stipulated block representation
\[
\mathcal A=\bigoplus_i\bigl(M_{d_i}(\C)\otimes I_{m_i}\bigr),\qquad d_i>0,
\]
the proof establishes that $k$ seeds generate all blocks simultaneously exactly when $m_i\le kd_i$ for every $i$. It follows that
\[
\boxed{\srank_{\mathcal A}(\mathcal H)=\max_i\left\lceil\frac{m_i}{d_i}\right\rceil.}
\]
The final Lean proposition is an \code{IsLeast} statement over the set of feasible seed counts. It proves both attainment and minimality. For an empty block family the finite supremum is defined as $0$, matching the zero module.





\section{Schmidt rank and the reduced state}\label{supp3:section:002}
Use coordinates $\psi_{ij}$ in the standard product basis and let $C_{ij}=\psi_{ij}$. An explicit linear equivalence identifies the state vector on the product index set with its coefficient matrix. The implementation checks the actual Kronecker action:
\[
\operatorname{coeff}\bigl((A\otimes I_B)\psi\bigr)=AC.
\]
It also proves that the transported matrix orbit equals the span of the state vectors obtained by this Kronecker action. The dimension formula therefore applies to the physical local action on the state space.

\subsection{An independent definition of Schmidt rank}
The source defines Schmidt rank as the least number of summands that are product vectors. A decomposition of length $s$ is the concrete equality
\[
C_{ij}=\sum_{a=1}^{s}u_a(i)v_a(j).
\]
It factors $C=UV$ through an $s$-dimensional intermediate space, so $\rank C\le s$. Conversely, choose a basis of the row space of $C$. Expanding each row in this basis constructs a decomposition with exactly $\rank C$ product terms. Thus
\[
\boxed{r(\psi)=\rank C.}
\]
The definition uses a least natural number, and the equality with matrix rank is a theorem. This argument establishes the minimum number of product terms. It does not construct an orthonormal Schmidt decomposition or prove a new theorem on singular value decomposition.

The matrix-orbit dimension and the coefficient equivalence now yield
\[
\boxed{\dim_\C\spanof\{(A\otimes I_B)\psi:A\in M_{d_A}(\C)\}
=d_A r(\psi).}
\]

\subsection{The actual partial trace and its support}
The reduced matrix is defined by tracing the addressed index of the density matrix of the pure state:
\[
(\rho_B)_{j\ell}=\sum_i\psi_{ij}\overline{\psi_{i\ell}}.
\]
For the coefficient convention used here, the exact Gram identity is
\[
\rho_B=C^{\mathsf T}(C^{\mathsf T})^*,
\qquad \rank\rho_B=\rank C=r(\psi).
\]
The transpose is important. With this convention, writing $C^*C$ instead would conjugate the reduced matrix. The Lean source checks the entrywise identity and positivity of the resulting matrix.

Full support is expressed as surjectivity of the reduced matrix's linear map. Assuming $d_A>0$, the checked equivalence is
\[
\boxed{\mathcal O_A(\psi)=\mathcal H_A\otimes\mathcal H_B
\quad\Longleftrightarrow\quad \Ran\rho_B=\mathcal H_B.}
\]
Normalization of $\psi$ is unnecessary for these dimension and support statements. The zero state has rank $0$ because its decomposition needs no product terms.





\section{Constructed cyclic character projectors}\label{supp3:section:003}
Let $N>0$ and let $S$ be a unitary complex matrix with $S^N=I$. The implementation uses the standard additive character of $\mathbb Z/N\mathbb Z$ and proves its expression through a root of unity:
\[
\omega=\exp(2\pi i/N),\qquad
\chi(r)=\omega^{\operatorname{val}(r)}.
\]
Here $\operatorname{val}(r)$ is the canonical representative in $\{0,\ldots,N-1\}$. The representation is constructed as
\[
\rho(r)=S^{\operatorname{val}(r)}.
\]
Reduction of exponents modulo $N$ is proved from $S^N=I$, giving $\rho(r+s)=\rho(r)\rho(s)$ and $\rho(0)=I$. The final matrix theorem derives the representation law and character orthogonality from the stated inputs.

\subsection{Exact coefficient formula}
The packet definition is a finite sum over residue classes:
\[
P_j=\frac1N\sum_{r\in\mathbb Z/N\mathbb Z}\chi(-rj)\rho(r).
\]
The source proves
\[
\chi(-rj)=\bigl(\omega^{\operatorname{val}(r)\operatorname{val}(j)}\bigr)^{-1},
\]
so this is precisely the manuscript's $N$-term Fourier projector formula, with canonical representatives.

\subsection{Algebraic identities}
Primitivity of the standard character yields the checked finite sum
\[
\sum_r\chi(tr)=\begin{cases}N,&t=0,\\0,&t\ne0.\end{cases}
\]
Reindexing finite sums then gives
\[
\boxed{P_jP_k=\delta_{jk}P_j,\qquad \sum_jP_j=I,\qquad SP_j=\omega^jP_j.}
\]
Multiplicativity and completeness are first proved for a general representation in a complex algebra, then instantiated with the actual powers of $S$. This organization makes the algebra of finite sums reusable while preserving the final hypotheses.

\subsection{Orthogonality of the projectors}
Unitarity supplies $\rho(r)^*=\rho(-r)$, while the character obeys $\overline{\chi(r)}=\chi(-r)$. Reindexing $r\mapsto-r$ proves $P_j^*=P_j$. The final concrete matrix theorem includes Hermiticity together with multiplication, completeness, and eigenaction. Together these identities establish orthogonal spectral packets.

The proof includes the degenerate modulus $N=1$: finite order then forces $S=I$ and the sole projector is $I$. It does not need an artificial assumption $N>1$.





\section{The actual Euclidean transfer exponential}\label{supp3:section:004}
Let $H$ be a Hermitian matrix on a nonzero complex state space of finite dimension. The source defines
\[
T_e(t)=\exp\bigl(-t(H-eI)\bigr)
\]
using Lean's exponential defined by a power series. The proof then relates this defined operator to its spectral expansion.

\subsection{Connecting power series and spectral calculus}
The shifted generator is proved Hermitian. The implementation then identifies the actual matrix exponential with continuous functional calculus:
\[
\boxed{T_e(t)=f_t(H),\qquad f_t(x)=e^{-t(x-e)}.}
\]
The formal proof uses composition, scalar multiplication, subtraction, and the established agreement between the exponential in functional calculus and the exponential defined by a power series. Spectral mapping is consequently available for the operator that was actually defined:
\[
\Spec_\R T_e(t)=\{e^{-t(x-e)}:x\in\Spec_\R H\}.
\]
Hermiticity also identifies the full complex spectrum with the image of these real eigenvalues. A separate checked theorem states the interval containing the complex spectrum explicitly.

\subsection{The minimum eigenvalue}
The ground energy is the minimum of the finite, nonempty eigenvalue set. The proof establishes both
\[
E_0\in\Spec_\R H,\qquad E_0\le x\quad(x\in\Spec_\R H).
\]
The assumption that the state space is nonzero is recorded in the theorem because the minimum of an empty spectrum is not defined by the manuscript's formula.

\subsection{Positivity, semigroup composition, and the interval}
Every eigenvalue of $T_e(t)$ is strictly positive, for arbitrary real $e,t$, since the real exponential is positive. For $e=E_0$ and $t\ge0$ it is also at most one. Therefore
\[
\boxed{T_{E_0}(t)>0,\qquad \Spec T_{E_0}(t)\subseteq(0,1].}
\]
The two scalar multiples of $H-eI$ commute, so the exponential defined by the power series obeys
\[
\boxed{T_e(t+s)=T_e(t)T_e(s),\qquad T_e(0)=I.}
\]
The semigroup identity is proved for all real $t,s$. Nonnegative time is needed for the upper spectral bound, not for the composition identity or strict positivity.

These results concern operators on a finite dimensional space. Continuum transfer limits, gravitational path integrals, and physical spacetime models require further structure. The numerical figures retain their earlier computation records and were not recomputed during this prose revision.





